\newcommand*{\orcid}{}
\newcommand{\data}[1]{\textit{#1}}

% This command modifies the underline to accomodate drop characters
\renewcommand{\ULdepth}{1.8pt}
\contourlength{0.8pt}

\newcommand{\myuline}[1]{%
  \uline{\phantom{#1}}%
  \llap{\contour{white}{#1}}}%

% This command is used to mark inline data as ! or ? and italicize the data. Don't use for *.
\newcommand{\mymark}[2]{\textsuperscript{#1}\textit{#2}}

% This command is used to create an initial example with superscript ! or ? and italicize the data. Don't use for *.
\newcommand{\eamark}[2]{\ea[\textsuperscript{#1}]{\textit{#2}}}

% This command is used to create a subsequent example with superscript ! or ? and italicize the data. Don't use for *.
\newcommand{\exmark}[2]{\ex[\textsuperscript{#1}]{\textit{#2}}}

% Define a command for strikethrough text
\newcommand{\textst}[1]{
  \tikz[baseline=(text.base)]{
    \node[inner sep=0pt,outer sep=0pt] (text) {#1};
    \draw[line width=0.5pt] ([yshift=1ex]text.south west) -- ([yshift=1ex]text.south east); % Adjust yshift for line position
  }
}

%----------------------------------------------------------------------
% Node labels in CGEL trees are defined with \Node, 
% which is defined so that \Node{Abcd}{Xyz} yields 
% a label with the function Abcd on the top, in small
% sanserif font, followed by a colon, and the category 
% Xyz on the bottom.
\newcommand{\Node}[2]{\small\textsf{#1:}\\{#2}}
% For commonly used functions this is defined with \(function)
\newcommand{\Head}[1]{\Node{Head}{#1}}
\newcommand{\Subj}[1]{\Node{Subj}{#1}}
\newcommand{\Comp}[1]{\Node{Comp}{#1}}
\newcommand{\Mod}[1]{\Node{Mod}{#1}}
\newcommand{\Det}[1]{\Node{Det}{#1}}
\newcommand{\PredComp}[1]{\Node{PredComp}{#1}}
\newcommand{\Crd}[1]{\Node{Coordinate}{#1}}
\newcommand{\Mk}[1]{\Node{Marker}{#1}}
\newcommand{\Obj}[1]{\Node{Obj}{#1}}
\newcommand{\Sup}[1]{\Node{Supplement}{#1}}